\documentclass[12pt]{article}
\usepackage[margin=1in]{geometry}
\usepackage{amsmath,amssymb,amsthm}
\usepackage{natbib}
\usepackage{hyperref}
\usepackage{booktabs}
\usepackage{graphicx}
\usepackage{setspace}
\usepackage{enumerate}

\hypersetup{colorlinks=true,linkcolor=blue,citecolor=blue,urlcolor=blue}

\newtheorem{theorem}{Theorem}
\newtheorem{proposition}{Proposition}
\newtheorem{corollary}{Corollary}
\newtheorem{definition}{Definition}

\title{Zero-Collapse and Resource Allocation in Multiplicative Systems: \\
Evidence from Housing, Austerity, and Development}
\author{David Kirsch\thanks{Independent researcher. Email: \texttt{david@davidkirsch.me}. Web: \url{https://davidkirsch.me}.}}
\date{July 2026}

\begin{document}
\maketitle

\begin{abstract}
\noindent This paper applies the O-Ring Inversion --- the mathematical result that marginal returns to investment are steepest at the weakest dimension in multiplicative systems --- to three domains of resource allocation: homelessness policy, fiscal crisis response, and cross-country development. In Paper~1 of this research program (Kirsch, 2026a), we proved that when independent, non-substitutable dimensions compose an emergent property, the unique functional form satisfying five structural axioms is the weighted geometric mean, $f(\mathbf{x}) = k\prod x_i^{\alpha_i}$, and that the partial derivative $\partial f/\partial x_i = (f/x_i)\alpha_i$ is maximized when $x_i$ is minimized. Here we test three predictions of that result against observational and quasi-experimental data. First, Housing First programs --- which address the zero-valued dimension (housing stability) before other interventions --- produce cross-dimensional gains consistent with the multiplicative prediction, achieving 71\% stable housing versus 29\% in controls and reducing psychiatric emergency use by 38\%. Second, the Greece--Iceland natural experiment (2008--2015) shows that fiscal austerity cutting dimensions near zero produces GDP contraction (26\%) exceeding additive model predictions, while Iceland's floor-preserving response yielded faster recovery. Third, in a 162-country cross-section, governance quality interacts multiplicatively with other development dimensions: countries with governance near zero show near-zero returns to capital investment regardless of human capital or infrastructure endowment. These results are consistent with the binding constraint principle of Hausmann, Rodrik \& Velasco (2005) and extend Kremer's (1993) O-Ring production function from task-level complementarity to system-level resource allocation. We identify conditions under which the multiplicative model fails and argue that the convergence of compassionate and efficient allocation is not a values claim but a consequence of the calculus.

\medskip
\noindent \textbf{Keywords:} multiplicative composition, O-Ring theory, zero-collapse, binding constraints, Housing First, austerity, development economics, resource allocation

\medskip
\noindent \textbf{JEL Classification:} O11, I38, H12, C23
\end{abstract}

\onehalfspacing

\section{Introduction}
\label{sec:intro}

The question of where to allocate scarce resources --- across dimensions of a system, among populations, between competing priorities --- is among the oldest in economics. Standard approaches typically assume additive separability: investments in different dimensions contribute independently to total output, and weakness in one dimension can be compensated by strength in another. This paper argues, on mathematical and empirical grounds, that this assumption is systematically wrong for a wide class of systems, and that the error has measurable consequences for policy.

The mathematical foundation is established in the companion papers of this research program. Paper~1 \citep{kirsch2026mc} proves that when an emergent property depends on multiple independent, non-substitutable dimensions, the unique functional form satisfying five structural axioms (dimensionwise monotonicity, independence, scale invariance, zero-collapse, and universality) is the weighted geometric mean:
\begin{equation}
\label{eq:mc}
f(\mathbf{x}) = k \prod_{i=1}^{n} x_i^{\alpha_i}, \quad \alpha_i > 0, \quad k > 0.
\end{equation}
Paper~2 \citep{kirsch2026geom} demonstrates that this functional form arises as the Riemannian volume density on the product statistical manifold equipped with the Fisher--Rao metric, establishing the axioms as geometric consequences rather than independent postulates.

The result that matters for resource allocation is the \emph{O-Ring Inversion}. The partial derivative of (\ref{eq:mc}) with respect to dimension $x_i$ is:
\begin{equation}
\label{eq:inversion}
\frac{\partial f}{\partial x_i} = \frac{f}{x_i} \cdot \alpha_i.
\end{equation}
Because this expression is inversely proportional to $x_i$, marginal returns are steepest at the weakest dimension. The optimal allocation of a marginal unit of investment is always to the dimension closest to zero --- not because this is compassionate, but because the partial derivative is largest there. Zero in any dimension collapses the entire product regardless of the magnitude of other dimensions. These are mathematical properties of the multiplicative form, not empirical claims.

This paper tests whether these mathematical properties predict real-world outcomes in three domains where resource allocation decisions have been made under implicit or explicit additive assumptions:

\begin{enumerate}
\item \textbf{Individual zero-collapse:} Housing First programs, which address the zero-valued dimension (housing stability) before attempting other interventions, versus treatment-as-usual programs that distribute resources additively across dimensions.
\item \textbf{Government crisis response:} The Greece--Iceland natural experiment, where two countries facing similar financial shocks adopted opposite policies --- Greece cutting dimensions near zero (social safety nets, healthcare, pensions) and Iceland preserving them.
\item \textbf{Cross-country development:} The interaction between governance quality and other development dimensions across 162 countries, testing whether near-zero governance annihilates returns to investment in human capital and infrastructure.
\end{enumerate}

The paper connects to two established literatures. Kremer's (1993) O-Ring theory of economic development proposes a multiplicative production function for task-level complementarity: output equals the product of worker quality across tasks, so a single low-quality worker can destroy the value of an entire production chain. Our contribution extends this from within-firm task complementarity to system-level resource allocation across policy dimensions. Hausmann, Rodrik \& Velasco's (2005) growth diagnostics framework argues that reformers should identify and target the most \emph{binding constraint} on growth, rather than pursuing broad-based reform. Our result provides a formal basis for this intuition: in a multiplicative system, the binding constraint is precisely the dimension where the partial derivative is largest, and addressing it produces the greatest total-system return.

We stress from the outset that the multiplicative model is not universally correct. It applies when dimensions are genuinely independent and non-substitutable --- when each is categorically necessary for the emergent property and cannot be replaced by abundance in another dimension. We identify specific conditions under which the model fails and report cases where additive specifications outperform.


\section{The Binding Constraint Principle}
\label{sec:theory}

\subsection{Mathematical Statement}

Consider a system whose output $E$ depends on $n$ independent dimensions $x_1, x_2, \ldots, x_n$, each measured on $[0, \infty)$. The five axioms of multiplicative composition (proved in \citet{kirsch2026mc}) yield the unique form (\ref{eq:mc}). The two properties relevant to resource allocation follow immediately.

\begin{proposition}[Zero-Collapse]
\label{prop:zero}
If $x_j = 0$ for any $j \in \{1, \ldots, n\}$, then $f(\mathbf{x}) = 0$ regardless of the values of $x_i$ for $i \neq j$.
\end{proposition}

\begin{proposition}[O-Ring Inversion]
\label{prop:inversion}
For any two dimensions $i, j$ with $\alpha_i = \alpha_j$, if $x_i < x_j$ then $\partial f / \partial x_i > \partial f / \partial x_j$. More generally, the marginal return to investment in dimension $i$ is proportional to $\alpha_i / x_i$.
\end{proposition}

The first property means that no amount of investment in strong dimensions can compensate for a dimension at zero. The second means that the optimal allocation of a marginal unit of investment is to the dimension with the highest ratio $\alpha_i / x_i$ --- typically the weakest dimension, unless dimension weights differ dramatically.

\subsection{Quantitative Illustration}

Consider a five-dimensional system with equal weights $\alpha_i = 0.2$ and current values $(0.8, 0.7, 0.6, 0.5, 0.1)$. The current output is $f = 0.8^{0.2} \times 0.7^{0.2} \times 0.6^{0.2} \times 0.5^{0.2} \times 0.1^{0.2} \approx 0.437$. The marginal return to a 0.1-unit investment in each dimension is:

\begin{center}
\begin{tabular}{lcccc}
\toprule
Dimension & Current value & New value & New $f$ & $\Delta f$ \\
\midrule
$x_1$ & 0.8 & 0.9 & 0.447 & 0.010 \\
$x_2$ & 0.7 & 0.8 & 0.450 & 0.013 \\
$x_3$ & 0.6 & 0.7 & 0.454 & 0.017 \\
$x_4$ & 0.5 & 0.6 & 0.460 & 0.023 \\
$x_5$ & 0.1 & 0.2 & 0.538 & 0.101 \\
\bottomrule
\end{tabular}
\end{center}

The same 0.1-unit investment produces a return $10.1\times$ larger when directed at the weakest dimension than at the strongest. This is a general property of the multiplicative form.

\subsection{Relationship to Prior Work}

Kremer (1993) proposes the production function $y = n k B \prod_{i=1}^{n} q_i$, where $q_i$ is the quality of worker $i$ executing task $i$. This generates positive assortative matching: high-skill workers are matched together because the multiplicative form means each worker's marginal product depends on the quality of all other workers. Our framework generalizes Kremer's insight from within-firm task complementarity to system-level composition of independent dimensions, and from worker quality to policy-relevant variables (housing stability, governance quality, healthcare access).

Hausmann, Rodrik \& Velasco (2005) argue that economic growth is constrained by binding constraints and that reform should target the most binding one. Their growth diagnostics decision tree identifies whether the binding constraint is cost of finance, low returns to economic activity, poor governance, or inadequate infrastructure. Our Proposition~\ref{prop:inversion} provides a formal basis for this intuition: in a multiplicative system, the binding constraint is the dimension with the lowest $x_i$ (for equal weights), and addressing it produces the greatest return because $\partial f / \partial x_i \propto 1/x_i$.

The distinction from additive models is sharp. Under $f(\mathbf{x}) = \sum \beta_i x_i$, marginal returns are constant: $\partial f / \partial x_i = \beta_i$ regardless of $x_i$. The optimal allocation under additive composition is always to the dimension with the highest coefficient $\beta_i$, irrespective of current levels. Under multiplicative composition, optimal allocation depends on \emph{both} the weight $\alpha_i$ and the current level $x_i$, shifting dynamically toward the weakest dimension.


\section{Housing First and Individual Zero-Collapse}
\label{sec:housing}

\subsection{The Multiplicative Prediction}

If an individual's capacity to thrive composes multiplicatively from housing stability $\times$ healthcare access $\times$ employment readiness $\times$ social support $\times$ financial security, then a person with housing stability at zero receives zero return from investments in any other dimension. The multiplicative model predicts that addressing the zero dimension first should produce \emph{cross-dimensional} gains --- improvements not only in housing but simultaneously in health, employment, and social outcomes --- because the product is uncollapsed.

The standard alternative is the ``treatment first'' or ``continuum of care'' model, which distributes resources across dimensions additively. Under additive assumptions, each intervention adds value independently regardless of housing status. Under multiplicative assumptions, these investments produce near-zero returns until housing is addressed.

\subsection{Evidence}

The most rigorous evidence comes from the Canadian \emph{At Home / Chez Soi} randomized controlled trial, which randomized 2,148 participants with mental illness across five cities to Housing First with Assertive Community Treatment (ACT) or Intensive Case Management (ICM) versus treatment-as-usual (TAU). High-needs participants receiving Housing First with ACT spent 71\% of the study period in stable housing versus 29\% for TAU controls (adjusted absolute difference = 42\%, $p < 0.01$), entered housing in 73 days versus 220 days ($p < 0.001$), and reported higher quality of life (ASMD = 0.15, $p < 0.01$) and better community functioning (ASMD = 0.18, $p < 0.01$) over two years \citep{aubry2016oneyear, goering2014ahs}.

A Campbell Collaboration systematic review of 43 randomized controlled trials found that Housing First reduces the number of days spent homeless by approximately 62.5 days (95\% CI: $-86.86$ to $-38.14$) and increases days in stable housing by 110.1 days (95\% CI: 93.05 to 127.15) relative to usual services, with moderate certainty evidence \citep{munthekaas2018}.

A randomized trial of permanent supportive housing for chronically homeless high users of public services in Santa Clara County (N = 423) found that 86\% of those randomized to PSH received housing and remained housed for 92.9\% of the follow-up period. The intervention reduced psychiatric ED visits by 38\% (IRR = 0.62, 95\% CI: 0.43--0.91), shelter days by 70\% (IRR = 0.30, 95\% CI: 0.17--0.53), and increased outpatient mental health service use by 84\% (IRR = 1.84, 95\% CI: 1.43--2.37) \citep{raven2020}. The cross-dimensional spillover is consistent with the multiplicative prediction.

A VA demonstration study of Housing First for chronically homeless veterans found that the approach reduced time to housing placement from 223 to 35 days, with housing retention at 12 months of 98\% for Housing First versus 86\% for TAU (OR = 8.33, $p = 0.023$) \citep{montgomery2013}.

Finland's national Housing First program (2008--2022) reduced long-term homelessness by 68\%, with shelter beds in Helsinki declining from approximately 600 to near zero. Finland is the only EU country where homelessness decreased during this period.

\subsection{Multiplicative Interpretation}

The pattern across these studies is consistent with the multiplicative model. Housing First does not merely produce housing outcomes; it produces outcomes across dimensions that were not directly targeted. When housing stability moves from near-zero to nonzero, the product is uncollapsed, and investments that were previously multiplying by zero begin producing returns. The cross-dimensional spillover --- better mental health outcomes, reduced emergency service use, increased engagement with treatment --- is the signature of a multiplicative system in which the binding constraint has been relieved.

\subsection{Limitations}

Housing First studies do not directly test the multiplicative functional form. The evidence is consistent with the multiplicative prediction but does not rule out alternative explanations (reduced cognitive load, stable address for service access). Furthermore, Housing First does not improve substance use outcomes relative to TAU \citep{somers2015}, suggesting that the cross-dimensional effect is not universal. The multiplicative model predicts that other investments \emph{can} produce returns once the zero is addressed, not that they will automatically do so without investment.


\section{Austerity, Stimulus, and the Greece--Iceland Natural Experiment}
\label{sec:austerity}

\subsection{The Multiplicative Prediction}

If a nation's economic output composes multiplicatively from institutional stability $\times$ human capital $\times$ infrastructure $\times$ social safety nets $\times$ financial system health, then cutting the dimensions closest to zero during a crisis should produce disproportionate output decline --- cliff-like contraction, not proportional decline. Conversely, investing in the weakest dimensions should produce disproportionate recovery.

\subsection{The Natural Experiment}

The 2008 global financial crisis created a natural experiment. Greece and Iceland faced similar triggers --- financial system collapse, sovereign debt distress --- but adopted opposite responses.

\textbf{Greece (austerity, 2010--2015):} Three bailout programs totaling \texteuro289 billion imposed deep fiscal consolidation: cuts to healthcare spending, pensions, public sector employment, and social safety nets. The dimensions closest to zero were cut furthest.

\textbf{Iceland (floor-preserving, 2008--2012):} Iceland allowed its banking system to fail, devalued its currency, and maintained social safety nets. The three largest banks were placed in receivership; their foreign obligations were not socialized.

\subsection{Outcomes}

Greek GDP declined by 26--27\% from peak to trough over 2008--2013, with unemployment reaching 28\% and youth unemployment 58\%. The contraction exceeded the U.S.\ Great Depression in duration. Critically, the austerity program failed on its own terms: debt-to-GDP rose from approximately 130\% (2009) to over 180\% (2014). Health inequalities widened: unmet medical need due to cost doubled from 7\% to 13.9\% between 2008 and 2013 \citep{economou2016health}.

Iceland's GDP returned to pre-crisis levels by 2015. Unemployment fell to approximately 3\% by 2015. IMF loans were repaid early (2012). GDP growth reached 4.1\% by 2015.

\subsection{The Blanchard--Leigh Multiplier Revision}

The IMF's own analysis provides independent evidence. \citet{blanchard2013} found that the fiscal multiplier during the crisis was approximately 1.5, substantially higher than the 0.5 implicitly assumed in projections. Each euro of spending cuts reduced GDP by approximately three times more than models predicted. The systematic underestimation is consistent with the multiplicative model: cutting dimensions near zero produces cliff-like contraction that linear models underpredict.

\subsection{Difference-in-Differences Framework}

A formal specification treats the divergent responses as the treatment:
\begin{equation}
\label{eq:did}
\ln Y_{ct} = \beta_0 + \beta_1 \text{Post}_t + \beta_2 \text{Austerity}_c + \beta_3 (\text{Post}_t \times \text{Austerity}_c) + \gamma \mathbf{X}_{ct} + \epsilon_{ct}
\end{equation}
testing whether $\beta_3 < 0$. Covariates include trade openness, financial sector exposure, and eurozone membership. The two-country design lacks statistical power; the natural experiment is suggestive, not definitive. The direction and magnitude of the divergence are consistent with the multiplicative prediction.


\section{Cross-Country Development: The Governance Zero}
\label{sec:development}

\subsection{The Multiplicative Prediction}

If national output composes multiplicatively from governance $\times$ human capital $\times$ infrastructure $\times$ health, then countries with governance near zero should exhibit near-zero output regardless of other endowments.

\subsection{Cross-Section Analysis}

Using World Bank WDI for 162 countries, we estimate additive and multiplicative specifications:

\noindent Additive:
\begin{equation}
\ln Y_c = \beta_0 + \beta_1 G_c + \beta_2 H_c + \beta_3 I_c + \beta_4 S_c + \epsilon_c
\end{equation}

\noindent Multiplicative (log-linearized):
\begin{equation}
\ln Y_c = \alpha_0 + \alpha_1 \ln G_c + \alpha_2 \ln H_c + \alpha_3 \ln I_c + \alpha_4 \ln S_c + \epsilon_c
\end{equation}

In prior analysis \citep{kirsch2026mc}, the multiplicative specification achieves a Spearman rank correlation of 0.953 versus 0.726 for the additive specification (1.31$\times$ advantage). Countries with governance near zero --- Burundi, South Sudan, Somalia --- exhibit GDP far below additive predictions but consistent with multiplicative zero-collapse.

\subsection{Panel Extension}

The panel specification with fixed effects:
\begin{equation}
\ln Y_{ct} = \mu_c + \lambda_t + \delta_1 \ln G_{ct} + \delta_2 \ln H_{ct} + \delta_3 \ln I_{ct} + \delta_4 \ln S_{ct} + \delta_5 (\ln G_{ct} \times \ln H_{ct}) + \nu_{ct}
\end{equation}
tests whether $\delta_5 > 0$ after absorbing country and year fixed effects. IV estimation using colonial origin instruments \citep{acemoglu2001colonial} can provide additional identification.

\subsection{Failed State Trajectories}

\textbf{Venezuela:} GDP fell from \$372.6 billion (2012) to approximately \$82 billion (2024), an effective contraction of 78\%. Proven oil reserves remained among the world's largest (303 billion barrels). WGI Rule of Law declined from $-0.76$ (1998) to $-2.35$ (2020), reaching the 0th percentile. Resource endowment unchanged; governance collapsed and output collapsed disproportionately.

\textbf{Zimbabwe:} Following land reform (2000), GDP per capita declined from \$1,640 to \$661 (1998--2008), hyperinflation peaked at 79.6 billion percent monthly (November 2008). WGI Rule of Law reached the 1st percentile.

In both cases, additive models predict 16--20\% decline if governance (one of four dimensions) fails. Actual declines of 60--88\% are consistent with multiplicative cliff-collapse.


\section{When the Multiplicative Model Fails}
\label{sec:limitations}

\subsection{High Substitutability}

When dimensions are substitutable, the additive model outperforms. In Paper~1, gender equity showed additive-dominant composition (Spearman ratio $0.89\times$), consistent with partial substitutability among gender equity dimensions.

\subsection{Threshold Effects}

The multiplicative model predicts smooth collapse as dimensions approach zero. Real systems sometimes exhibit discrete threshold effects that the continuous multiplicative form approximates but does not capture exactly.

\subsection{Education Ceiling Effects}

Education quality across 40 countries showed near-parity between multiplicative and additive specifications ($0.97\times$), suggesting significant substitutability among educational dimensions.

\subsection{Measurement Limitations}

WGI governance indicators are perception-based composites. Year-to-year changes are frequently not statistically significant, complicating panel analysis. CES estimation of the substitution parameter is sensitive to specification in small samples.

\subsection{Scope Condition}

The multiplicative law applies when: (a) the emergent property depends on each dimension (AND-gated), and (b) dimensions are non-substitutable. Five missing fibers do not destroy a rug because fibers are substitutable. But zero press freedom produces zero democracy because press freedom is categorically different from electoral integrity.


\section{Discussion}
\label{sec:discussion}

\subsection{The Institutional Additive Default}

Institutions default to additive resource allocation. Treatment-as-usual programs distribute services additively. Austerity cuts proportionally across dimensions. The Washington Consensus prescribed broad-based reform without prioritizing binding constraints. In each case, the multiplicative model predicts suboptimal outcomes when any dimension is near zero, and the evidence confirms.

The binding constraint principle of \citet{hrv2005} is the clearest policy articulation of the multiplicative insight. Our contribution formalizes it: the binding constraint is a mathematical consequence of the multiplicative form.

\subsection{Compassionate Allocation as Efficient Allocation}

The O-Ring Inversion produces a result with ethical implications: in multiplicative systems, the investment maximizing total output targets the weakest dimension. When the weakest dimension affects the most vulnerable --- housing for the homeless, safety nets during crises, governance in failed states --- the compassionate and efficient allocations are mathematically identical.

This is consistent with \citet{ostry2014redistribution}, who find redistribution generally pro-growth or growth-neutral; with \citet{egger2022kenya}, who found 2.5--2.7$\times$ fiscal multipliers for transfers to poor households; and with \citet{banerjee2015brac}, who found that multidimensional ``big push'' interventions produced 133--433\% returns, outperforming single-dimension interventions within the same RCT.

\subsection{Relationship to the O-Ring Literature}

Kremer (1993) established that multiplicative production functions generate positive assortative matching and large income differences. Jones (2013) extended this with a foolproof sector. Our contribution extends the O-Ring framework from production theory to resource allocation: the same multiplicative structure that generates assortative matching within firms generates binding constraint dynamics across policy dimensions.


\section{Conclusion}
\label{sec:conclusion}

This paper has tested three predictions of the O-Ring Inversion against observational and quasi-experimental data from housing policy, fiscal crisis response, and cross-country development.

The evidence is consistent with the predictions. Housing First programs that address the zero-valued dimension produce cross-dimensional gains that treatment-as-usual does not. Austerity that cuts dimensions near zero produces cliff-like contraction that additive models underpredict. Countries with governance near zero show near-zero returns to capital investment regardless of other endowments.

These results do not prove that all social and economic systems compose multiplicatively. They show that the multiplicative model produces testable predictions that the additive model does not, that these predictions are confirmed in the domains examined, and that the additive default produces measurably suboptimal outcomes when any dimension is near zero.

The implications for resource allocation are direct. In systems where the multiplicative model applies, the optimal allocation is always to the weakest dimension. The compassionate allocation is the efficient allocation. The returns are at the bottom.

This is not an argument for any particular political position. It is a mathematical result with policy consequences. The partial derivative does not have a political opinion. It has a sign and a magnitude, and both point to the same place: the dimension closest to zero.


\begin{thebibliography}{99}

\bibitem[Acemoglu et al.(2001)]{acemoglu2001colonial}
Acemoglu, D., Johnson, S., \& Robinson, J.~A. (2001).
The colonial origins of comparative development: An empirical investigation.
\emph{American Economic Review}, 91(5), 1369--1401.

\bibitem[Aubry et al.(2016)]{aubry2016oneyear}
Aubry, T., Goering, P., Veldhuizen, S., et al. (2016).
A multiple-city RCT of Housing First with Assertive Community Treatment for homeless Canadians with serious mental illness.
\emph{Psychiatric Services}, 67(3), 275--281.

\bibitem[Banerjee et al.(2015)]{banerjee2015brac}
Banerjee, A., Duflo, E., Goldberg, N., et al. (2015).
A multifaceted program causes lasting progress for the very poor: Evidence from six countries.
\emph{Science}, 348(6236), 1260799.

\bibitem[Blanchard \& Leigh(2013)]{blanchard2013}
Blanchard, O.~J., \& Leigh, D. (2013).
Growth forecast errors and fiscal multipliers.
\emph{American Economic Review}, 103(3), 117--120.

\bibitem[Economou et al.(2016)]{economou2016health}
Economou, C., Kaitelidou, D., Kentikelenis, A., et al. (2016).
The impact of the financial crisis on the health system and health in Greece.
In \emph{Economic Crisis, Health Systems and Health in Europe}. WHO/European Observatory.

\bibitem[Egger et al.(2022)]{egger2022kenya}
Egger, D., Haushofer, J., Miguel, E., et al. (2022).
General equilibrium effects of cash transfers: Experimental evidence from Kenya.
\emph{Econometrica}, 90(6), 2603--2643.

\bibitem[Goering et al.(2014)]{goering2014ahs}
Goering, P., Veldhuizen, S., Watson, A., et al. (2014).
National At Home/Chez Soi final report.
Mental Health Commission of Canada.

\bibitem[Hausmann, Rodrik, \& Velasco(2005)]{hrv2005}
Hausmann, R., Rodrik, D., \& Velasco, A. (2005).
Growth diagnostics.
Working paper, Harvard University.

\bibitem[Jones(2013)]{jones2013oring}
Jones, G. (2013).
The O-Ring sector and the foolproof sector.
\emph{Journal of Economic Behavior \& Organization}, 87, 183--199.

\bibitem[Kirsch(2026a)]{kirsch2026mc}
Kirsch, D. (2026a).
Multiplicative composition of independent dimensions in emergent systems.
Working paper. \url{https://davidkirsch.me/papers/multiplicative-composition.pdf}.

\bibitem[Kirsch(2026b)]{kirsch2026geom}
Kirsch, D. (2026b).
Information geometry of multiplicative composition.
Working paper.

\bibitem[Kremer(1993)]{kremer1993oring}
Kremer, M. (1993).
The O-Ring theory of economic development.
\emph{Quarterly Journal of Economics}, 108(3), 551--575.

\bibitem[Montgomery et al.(2013)]{montgomery2013}
Montgomery, A.~E., Hill, L.~L., Kane, V., \& Culhane, D.~P. (2013).
Housing chronically homeless veterans: Evaluating the efficacy of a Housing First approach to HUD-VASH.
\emph{Journal of Community Psychology}, 41(4), 505--514.

\bibitem[Munthe-Kaas et al.(2018)]{munthekaas2018}
Munthe-Kaas, H.~M., Berg, R.~C., \& Blaasvaer, N. (2018).
Effectiveness of interventions to reduce homelessness: A systematic review and meta-analysis.
\emph{Campbell Systematic Reviews}, 14(1), 1--281.

\bibitem[Ostry et al.(2014)]{ostry2014redistribution}
Ostry, J.~D., Berg, A., \& Tsangarides, C.~G. (2014).
Redistribution, inequality, and growth.
IMF Staff Discussion Note SDN/14/02.

\bibitem[Raven et al.(2020)]{raven2020}
Raven, M.~C., Niedzwiecki, M.~J., \& Kushel, M. (2020).
A randomized trial of permanent supportive housing for chronically homeless persons with high use of publicly funded services.
\emph{Health Services Research}, 55(S2), 797--806.

\bibitem[Somers et al.(2015)]{somers2015}
Somers, J.~M., Moniruzzaman, A., \& Palepu, A. (2015).
Changes in daily substance use among people experiencing homelessness and mental illness: 24-month outcomes following randomization to Housing First or usual care.
\emph{Addiction}, 110(10), 1605--1614.

\end{thebibliography}

\end{document}
